\section{Related Work}
\label{sec:related}

\paragraph{RL and distillation for diffusion alignment.}
Reinforcement learning has been widely adapted to align diffusion text-to-image
models to reward signals~\cite{liu2025flow,zheng2025diffusionnft,schulman2017proximal}.
Flow-GRPO~\cite{liu2025flow} casts flow-matching sampling as an MDP and optimizes
a GRPO-style objective; DiffusionNFT~\cite{zheng2025diffusionnft} performs online
reinforcement on the forward process. These methods excel at single rewards but
struggle to serve several objectives jointly. DiffusionOPD~\cite{li2026diffusionopd}
reframes multi-task alignment as \emph{on-policy distillation}: specialized
teachers, trained once per task, supervise a single student along its own rollout
via a closed-form per-step KL, extending the OPD idea~\cite{thinkingmachines2026opd}
from tokens to continuous diffusion states. We take DiffusionOPD as our base
method, reproduce it faithfully, and study how the multiple teachers should be
combined.

\paragraph{Multi-objective and multi-task optimization.}
The limitations of uniform loss summation are classical. Linear scalarization
recovers only the convex hull of the Pareto front~\cite{miettinen1999nonlinear,boyd2004convex},
whereas weighted Tchebycheff (min-max) scalarization can reach the entire front,
including non-convex regions~\cite{miettinen1999nonlinear}. In deep multi-task
learning, MGDA~\cite{sener2018multi,desideri2012mgda} finds a common descent
direction with a Pareto-stationarity guarantee, PCGrad~\cite{yu2020gradient}
projects away conflicting gradient components, and CAGrad~\cite{liu2021conflict}
takes a conflict-averse step that reduces the worst-case local objective while
converging to a stationary point of the average loss. These methods generally
assume commensurable objectives, an assumption violated by generic MTL where
losses have incompatible units and scales. A central theoretical point of our
work (\S\ref{sec:theory}) is that multi-teacher OPD \emph{does} satisfy
commensurability---all per-teacher losses are the same velocity-error functional
in one shared isotropic output metric---so adaptive Tchebycheff scalarization is
particularly well-posed here. We adopt scalarization rather than gradient surgery
and are careful not to over-transfer output-space isotropy to parameter space.

\paragraph{Loss weighting in diffusion training.}
Timestep-dependent loss weighting is common in diffusion training; Min-SNR~\cite{hang2023efficient}
treats denoising as multi-task learning across noise levels and down-weights
low-noise (high-SNR) steps to balance the objective. We test an analogous
per-step reweighting inside OPD (\S\ref{sec:ablation}) and find it harmful,
because high-noise steps set global structure that OPD must preserve---a useful
negative result that clarifies why OPD's uniform per-step weighting is a good
default. Robust regression via the Huber loss~\cite{huber1964robust} motivates our
trust-region variant on the velocity residual, which we find neutral.

\paragraph{Evaluation.}
We evaluate with the same reward suite as the base method: GenEval~\cite{ghosh2023geneval}
and OCR (rule-based), and PickScore~\cite{kirstain2023pick},
ClipScore~\cite{hessel2021clipscore}, HPSv2.1~\cite{wu2023human},
Aesthetics~\cite{schuhmann2022aesthetics}, ImageReward~\cite{xu2023imagereward},
and UnifiedReward~\cite{wang2025unified} (model-based), on
DrawBench~\cite{saharia2022photorealistic} prompts, with SD3.5-Medium~\cite{esser2024scaling}
as the base model.
